\documentclass[10pt,letterpaper]{article}
% Shared formatting for Joshua Cowell resume variants.
\usepackage[left=0.50in,right=0.50in,top=0.48in,bottom=0.48in]{geometry}
\usepackage[T1]{fontenc}
\usepackage{mathptmx}
\usepackage{enumitem}
\usepackage{xcolor}
\usepackage{hyperref}
\usepackage{tabularx}
\usepackage{array}
\usepackage{setspace}

\definecolor{navyblue}{RGB}{0,40,85}
\definecolor{darkgray}{RGB}{50,50,50}
\definecolor{medgray}{RGB}{95,95,95}

\hypersetup{colorlinks=false,pdfborder={0 0 0}}
\pagenumbering{gobble}
\setlength{\parindent}{0pt}
\setlength{\parskip}{0pt}
\setstretch{1.0}
\setlist[itemize]{
  leftmargin=0.16in,
  topsep=1.5pt,
  itemsep=0.6pt,
  parsep=0pt,
  partopsep=0pt,
  label=\textcolor{navyblue}{\textbullet}
}

\newcommand{\resumeSection}[1]{%
  \vspace{5pt}%
  {\normalsize\bfseries\textcolor{navyblue}{\MakeUppercase{#1}}}\par
  \vspace{1pt}%
  {\color{navyblue}\hrule height 0.5pt}%
  \vspace{3pt}%
}

\newcommand{\resumeEntry}[3]{%
  \textbf{#1} \hfill \textcolor{medgray}{#2} \\
  \textit{\textcolor{darkgray}{#3}}\par
}

\newcommand{\companyEntry}[2]{%
  \textbf{#1} \hfill \textcolor{medgray}{#2}\par
}

\newcommand{\roleLine}[2]{%
  \textbf{#1} \hfill \textcolor{medgray}{#2}\par
}
\newcommand{\currentRoleLine}[2]{%
  {\small\textit{\textcolor{darkgray}{Current: #1}}} \hfill \textcolor{medgray}{#2}\par
}


\newcommand{\earlierRoles}{%
  {\footnotesize\textit{Earlier roles: Assistant Chief Engineer; Senior Structural Engineer; Fin Fluid-Dynamics Engineer}}\par
}

\newcommand{\projectEntry}[1]{\textbf{#1}\par}

\newcommand{\resumeHeader}{%
  \begin{center}
    {\fontsize{20}{22}\selectfont\bfseries\textcolor{navyblue}{JOSHUA COWELL}}

    \vspace{3pt}

    {\small\textbf{\textcolor{navyblue}{Active Top Secret / SCI Clearance}}
    \enspace\textcolor{medgray}{$|$}\enspace
    \textbf{\textcolor{navyblue}{Full-Scope Polygraph}}}

    \vspace{3pt}

    {\small\color{darkgray}
      Lynchburg, VA
      \enspace\textcolor{medgray}{$|$}\enspace
      \href{mailto:joshuacowell2005@gmail.com}{joshuacowell2005@gmail.com}
      \enspace\textcolor{medgray}{$|$}\enspace
      \href{https://joshuacowell.com}{joshuacowell.com}
    }
  \end{center}
}

\newcommand{\educationBlock}[1]{%
  \resumeSection{Education}
  \textbf{Bachelor of Science in Mechanical Engineering} \hfill \textcolor{medgray}{Expected May 2027} \\
  \textit{\textcolor{darkgray}{Liberty University, Lynchburg, VA}} \hfill \textcolor{medgray}{GPA: 3.85}

  \vspace{2pt}
  {\small\textit{Selected Coursework:} #1}\par
  \vspace{1pt}
  {\small\textit{Founding Member:} American Nuclear Society (ANS) student chapter (2025); AIAA student chapter (2026)}\par
  {\small\textit{Leadership:} Liberty Rocketry -- Chairman, Board of Advisors (Jul 2026 -- Present)}\par
}

\newcommand{\researchEntry}{%
  \resumeEntry{Lead Undergraduate Research Assistant}{Sep 2024 -- Present}{Liberty University}
  \begin{itemize}
    \item Lead a five-person team conducting split-Hopkinson pressure bar (SHPB) testing for dynamic material characterization
    \item Develop test procedures and fixtures; use PicoScope strain data to generate stress-strain curves and correlate high-speed imaging with SEM results
  \end{itemize}
}


\begin{document}
\resumeHeader
\educationBlock{Fluid Dynamics, Heat Transfer, Thermodynamics I \& II, Thermal-Fluids Laboratory, Differential Equations, Materials~Engineering}

\resumeSection{Engineering Experience}

\resumeEntry{Thermal Hydraulics Engineering Intern}{May 2025 -- Dec 2025}{Framatome Nuclear}
\begin{itemize}
  \item Developed and validated a Python predictive-wear model against documented component wear, reducing prediction error by 75\% versus the legacy model
  \item Built a Python/xarray algorithm processing six-dimensional datasets to calculate safe operating limits, reducing analysis time from 800+ hours to minutes
  \item Performed system-level thermal-hydraulic and flow-induced-vibration (FIV) analyses using AFT Fathom and custom tools to support component qualification
  \item Authored technical documentation and presented findings to engineering leadership; received a \$2,500 technical-presentation award
\end{itemize}

\vspace{2pt}
\resumeEntry{Advanced Weapons Systems Analyst Intern}{May 2026 -- Present}{Central Intelligence Agency (CIA)}
\begin{itemize}
  \item Conduct testing of warhead systems and analyze results to calibrate simulation models
  \item Develop a Python-based geospatial analysis tool for operational planning that meets team requirements
  \item Brief senior Agency officials and write concise analytic products that communicate complex technical data in decision-ready formats
\end{itemize}

\vspace{2pt}
\resumeEntry{Chief Systems Engineer and Team Lead}{Jun 2025 -- Jul 2026}{Liberty Rocketry}
\earlierRoles
\begin{itemize}
  \item Placed 18th of 141 international collegiate teams at the 2026 IREC
  \item Led a 60-member student team in the design, testing, and launch readiness of a high-powered rocket
  \item Oversaw engineering decisions across propulsion, avionics, recovery, and aerodynamics subsystems
  \item Facilitated communication among internal teams, external vendors, and university faculty to improve project execution
  \item Reviewed SolidWorks, ANSYS, MATLAB, OpenRocket, and RocketPy models and simulations used for team design decisions
\end{itemize}

\vspace{2pt}
\researchEntry

\resumeSection{Technical Skills}
\begin{tabularx}{\textwidth}{@{}>{\bfseries}l X@{}}
Thermal \& Analysis: & Thermal-hydraulics, heat transfer, FIV, ANSYS Fluent, ANSYS Mechanical, AFT Fathom, CFD, FEA \\
Programming: & Python (NumPy, pandas, xarray), MATLAB, scientific computing, data analysis, automation \\
CAD \& Manufacturing: & SolidWorks, Onshape, GD\&T, technical drawings, 3D printing, CNC machining \\
Laboratory: & SHPB testing, SEM, high-speed imaging, test fixtures, material testing
\end{tabularx}

\resumeSection{Selected Projects}
\projectEntry{Heat Exchanger Thermal-Hydraulic Optimization \textnormal{\textit{-- MATLAB, ANSYS Fluent, SolidWorks}}}
\begin{itemize}
  \item Wrote MATLAB optimization code that evaluated 104 million crossflow and 287 million annular geometries against thermal, dimensional, and annual-cost constraints
  \item Validated the selected annular design against ANSYS Fluent CFD across heat-transfer and effectiveness metrics
\end{itemize}

\projectEntry{3D-Printed Model Rocket \textnormal{\textit{-- ANSYS Fluent, OpenRocket, Additive Manufacturing}}}
\begin{itemize}
  \item Designed a reusable threaded PET-GF15/PETG airframe and hybrid recovery system; simulations predicted Mach 0.5+ and approximately 3,000 ft altitude
\end{itemize}

\end{document}
